\section{Introduction}\label{sec:introduction}

In a shape optimization problem for thermal insulation, one may choose
both the heated body and the region surrounding it. The question is
whether a spherical body with a concentric insulating layer uses a
given amount of material most efficiently. For a stationary model
with quadratic bulk energy, the boundary loss is described by a
nonnegative potential of the temperature. The choice of that
potential distinguishes convective and radiative boundary laws
\cite{BNNT}.

Bucur, Nahon, Nitsch and Trombetti \cite{BNNT} proved concentric-ball
optimality for the quadratic boundary potential. They also obtained
general results for a volume penalty and certain constrained cases
with inactive volume constraint. Their question after Remark~4 asks
whether the concentric comparison holds for more general boundary
potentials, including the radiative law in their equation~(3).
The penalized comparison does not by itself give the constrained
comparison at every prescribed volume. We prove that comparison
for the full radiative potential, with all its parameters retained.

Here is the precise problem. Fix an integer \(n\ge2\), numbers
\(\beta>0\), \(\gamma\ge0\), and \(R>1\), and set
\begin{equation}\label{eq:theta}
 \Theta(s)=\beta\left(\frac{s^5}{5}+\gamma s^4+
                                  2\gamma^2s^3+2\gamma^3s^2\right).
\end{equation}
For a bounded Lipschitz open set \(\Omega\subset\R^n\) and a measurable
set \(K\subset\Omega\), let
\begin{equation}\label{eq:energy}
 E_\Theta(K,\Omega)=
 \min_{\substack{v\in H^1(\Omega),\ 0\le v\le1\\
                         v=1\ {\rm a.e.\ on}\ K}}
 \left\{\int_\Omega|\nabla v|^2\dd x+
          \int_{\partial\Omega}\Theta(\Tr v)\dd\mathcal H^{n-1}\right\}.
\end{equation}
The interval bounds are understood almost everywhere, and \(\Tr v\)
is the interior trace. Write the expression in braces as
\(\mathcal E_\Omega(v)\). We use the unit ball \(B_1\) to normalize
the heated volume and write \(\omega_n=|B_1|\).

\begin{theorem}\label{thm:main}
For every choice of \(n\ge2,\beta>0,\gamma\ge0,R>1\) above, every bounded
Lipschitz open \(\Omega\subset\R^n\), and every measurable \(K\subset\Omega\)
satisfying
\[
 |K|=\omega_n,\qquad |\Omega|\le\omega_nR^n,
\]
one has
\begin{equation}\label{eq:main}
 E_\Theta(K,\Omega)\ge\min_{1\le r\le R}E_\Theta(B_1,B_r).
\end{equation}
At \(r=1\), the pair is \(K=\Omega=B_1\).
\end{theorem}

Here \(B_r\) denotes the ball of radius \(r\) with the same center as
\(B_1\), and the energy on the right is the same trace energy
\eqref{eq:energy}. No connectedness of \(\Omega\) or regularity of
\(K\) is required. The minimum includes the bare ball and can use
less than the full available outer volume. This is the radiative
case of the constrained concentric-ball question, with the
upper-volume convention used in \cite{BNNT}.
We make no assertion here for an arbitrary increasing boundary potential.

The proof uses the boundary potential itself as the rearranged
variable. Extending that potential by zero records the outer
boundary loss as variation. Isoperimetry then gives a radial lower
cost, but the radial profiles may have jumps and need not represent
admissible \(H^1\) temperatures. We approximate their cost by profiles
with finitely many jumps. For the radiative polynomial, multiplying
the potential on the entire tail after a jump has a concave total
cost. One endpoint of the allowed multiplication interval removes
the jump, or removes the whole tail, without raising the energy.
Repeating this operation leaves a genuine concentric-ball competitor.
The possibility of shortening the tail explains why smaller outer
radii, including the bare ball, belong in \eqref{eq:main}.

Section~\ref{sec:radial} gives the radial energy formula and the scalar
concavity property. Section~\ref{sec:rearrangement} proves the radial
lower comparison with all singular variation retained.
Section~\ref{sec:jumps} constructs the finite-jump approximation,
removes the jumps, and proves the theorem.
All functions are real, all volume integrals use Lebesgue measure,
and Hausdorff measure is normalized by
\(\mathcal H^{n-1}(\partial B_r)=n\omega_n r^{n-1}\).
